\documentclass[11pt]{article}
\usepackage[a4paper,margin=23mm]{geometry}
\usepackage[T1]{fontenc}
\usepackage{lmodern,amsmath,amssymb,amsthm,mathtools,microtype,booktabs}
\usepackage[colorlinks=true,linkcolor=blue,citecolor=blue,urlcolor=blue]{hyperref}
\newcommand{\EE}{\mathbb E}
\newcommand{\Om}{\Omega_O}
\newcommand{\ph}{\varphi}
\newcommand{\bn}[2]{\binom{#1}{#2}}
\title{The first joint Hamiltonian layers above a Slater state\\
\large Same-disorder mixed fourth moments in complex double-scaled SYK}
\author{Alan Day \and Chengheng\thanks{AI research assistant (OpenAI).}}
\date{7 November 2026\quad Local research draft}
\begin{document}\maketitle
\begin{abstract}
The Slater-adapted creation block of the original Hamiltonian has a
dressed bosonic vacuum law, but captures only part of the first
Hamiltonian excitation. We retain its correlations with the same
Hamiltonian and evaluate every joint word through degree four.
Finite-size Gaussian pairings reduce to elementary counts and the
known embedded-ensemble covariance channel. Their double-scaled limit
gives a positive seven-dimensional space of polynomial vectors of
degree at most two. The first Hamiltonian actions exhibit a new
commutator direction of squared norm $1-q$. The actual disorder
complement has normalized fourth moment $2+\sqrt q+2q$; replacing it
by another standard $q$-Gaussian mode would change the model. These
results determine the first layers, with the all-degree joint
representation and its bulk interpretation left as the next tasks.
\end{abstract}

\section{State, normalization, and the literature interface}
Use standard complex CAR, $\{c_i,c_j^\dagger\}=\delta_{ij}$, and
\begin{equation}
 H_N=\sqrt{w_N}\sum_{|A|=|B|=p}J_{AB}C_A^\dagger C_B,
 \quad w_N=\frac{4^pJ_s^2}{\lambda_N}\left(\frac{p!}{N^p}\right)^2,
 \quad\lambda_N=\frac{p^2}{N}.
\end{equation}
Here $C_A^\dagger$ has increasing site order, $C_A$ is its adjoint,
and the centered Hermitian Gaussian coefficients obey
$\EE J_{AB}J_{CD}=\delta_{AD}\delta_{BC}$. Overlapping $A,B$ are
included. Fix even $N$, $m=a=N/2$, an occupied set $O$ of size $m$,
its empty complement $E$, and $\Om=|O\rangle$. All moments use
\begin{equation}
 \ph_N(X)=\EE\langle\Om,X\Om\rangle.
\end{equation}
This is the annealed Slater vector state. It is not assumed stationary.
Take $p\to\infty$, $\lambda_N\to\lambda>0$, fixed $J_s>0$ and
fixed word length, with $\hbar=1$ and string-energy units for $H_N$.
The normalized $h_N$ below is dimensionless.

\newpage
The same-$J$ excitation block and normalizations are
\begin{align}
 R_N^\dagger&=\sum_{A\subset E,\ B\subset O}J_{AB}C_A^\dagger C_B,
 &L_N&=\bn mp\bn ap,\\
 D_N&=\bn mp\bn{a+p}p,
 &h_N&=\frac{H_N}{\sqrt{w_ND_N}},\qquad
 M_N^\dagger=\frac{R_N^\dagger}{\sqrt{L_N}}.
\end{align}
Every subset in the sum has size $p$; $M_N$ is the adjoint.
Set
\begin{equation}
 q=e^{-4\lambda},\quad s=\sqrt q,\quad\gamma=q^{1/4},\quad
 \gamma_N=\sqrt{L_N/D_N},\quad\delta=\sqrt{1-s}.
\end{equation}
The preceding stage proved the pure-$M$ fixed-word law
$M=q^{n/2}b$, $M^\dagger=b^\dagger q^{n/2}$ in its vacuum
state, and the first-step fraction $L_N/D_N\to s$.
Here we retain the rest of $h_N\Om$.

Rajgadia--Xu's state-adapted construction uses matter Gaussian
coefficients separate from those of $H$; see their HTML
Eqs.~(42)--(44) and the selection rule after (47) \cite{RX}.
That selection rule is not an input for the present same-$J$ model.
Already at degree two,
\begin{equation}\label{second}
 \ph_N(h_N^2)=\ph_N(M_NM_N^\dagger)=1,\qquad
 \ph_N(h_NM_N^\dagger)=\ph_N(M_Nh_N)=\gamma_N.
\end{equation}
All other two-letter moments vanish. Odd total degree vanishes by
Gaussian parity, but unequal numbers of $M$ and $M^\dagger$ need
not make a mixed word vanish.

\section{An exact finite-size table for the three pairings}
In this section only, use the raw alphabet $H,+,-$ for
$H_N/\sqrt{w_N},R_N^\dagger,R_N$. Omit subscripts on counts.
A Gaussian edge with two $H$ endpoints is type $G$; every other
edge is type $S$, restricted to the excitation block or its adjoint.
An $S$ edge is forward if its creation acts first (right to left).
An endpoint $+$ is a creation and $-$ is an annihilation, fixing the
opposite endpoint even when that endpoint is $H$.
Two equal non-$H$ endpoints are incompatible and give zero.

For four written slots $0,1,2,3$, the three pairings are
\begin{equation}
 (01)(23)\ \text{disjoint},\qquad
 (02)(13)\ \text{crossed},\qquad
 (03)(12)\ \text{nested}.
\end{equation}
In a nested pairing the first edge listed is outer. Define
\begin{align}
 r_N&=\frac{\bn{m-p}p\bn{a-p}p}{L},
 &T_N&=\bn{m-p}p\bn ap,\\
 W_k&=\bn mp\bn pk\bn ak,
 &U_N&=\sum_{k=0}^p W_k\bn{a-k}p\bn{m-k}p.
\end{align}
Binomials outside their natural range are zero. $K_N$ denotes
the raw crossed $HHHH$ contraction.

\newpage
\begin{center}
\begin{tabular}{llll}
\toprule
Edges / orientation & Disjoint & Crossed & Nested\\
\midrule
$G,G$ & $D^2$ & $K_N$ & $D^2$\\
$S,S$, both forward & $L^2$ & $L^2r_N$ & $L^2r_N$\\
$G,S$, $S$ forward & $DL$ & $LT_N$ & see below\\
$S$ outer, $G$ inner & --- & --- & $DL$\\
$G$ outer, $S$ inner forward & --- & --- & $U_N$\\
$S$ outer forward, $S$ inner backward & $0$ & $0$ & $L$\\
$G$ outer, $S$ inner backward & $0$ & $0$ & $L$\\
\bottomrule
\end{tabular}
\end{center}
All other backward cases are zero. Sum the three contractions,
then divide by $D^{n_H/2}L^{(4-n_H)/2}$ to obtain the normalized
moment. This specifies all $3^4$ ordered fourth words.

\subsection*{Derivation of the entries}
For a fixed excitation label $\alpha=(A,B)$, write
$X_\alpha=C_A^\dagger C_B$. Relative to $\Om$, this is a pure
creation monomial of even degree $2p$. Two active excitation labels
must be disjoint on both $O$ and $E$. Their even CAR monomials then
commute, and each creation/adjoint pair has amplitude $+1$.
Disjoint time intervals give $L^2$; overlapping intervals give
$L\bn{m-p}p\bn{a-p}p$. If an inner selected edge is backward,
it can annihilate an outer selected excitation only by taking the
same entire label. It then recreates it, giving exactly $L$.
All other backward placements meet the vacuum boundary or cannot
remove the required particles and holes.

For general labels $Y_{AB}=C_A^\dagger C_B$, the channel
\begin{equation}\label{channel}
 \Phi(X)=\sum_{A,B}Y_{AB}X Y_{AB}^\dagger
\end{equation}
has $\Phi(I)=DI$ on the $m$-particle sector. Thus adjacent $G$
contractions and $G$ nested inside a selected edge give the table's
$D$ factors. The established embedded-ensemble spectrum is
\begin{equation}\label{eigen}
 \Lambda^r(p)=\bn{m-r}p\bn{N-m+p-r}p
\end{equation}
on rank-$r$ irreducible tensors \cite[Eqs.~(17),(25)--(26)]{Kota}.
A disjoint $p$-hop $X_\alpha$ is a rank-$p$ tensor: choose its
creation and annihilation indices as opposite extremal sets to
obtain a highest-weight vector; site permutations generate every
such disjoint hop. Hence $\Phi(X_\alpha)=T_NX_\alpha$, giving
the crossed $GS$ value $LT_N$. This is an application of the
known channel, not a newly discovered covariance eigenvalue.

\newpage
For $G$ outer and $S$ inner, let $k$ be the number of particles
that the first general hop moves from $O$ to $E$. The exact number
of allowed raw labels is $W_k$. Indeed choose its $p$ annihilated
occupied sites, the $k$ particles left displaced among those $p$,
and their $k$ empty destinations; the other $p-k$ are restored.
After this hop there are $a-k$ empty $E$ sites and $m-k$ filled
$O$ sites. A forward selected pair therefore has
$\bn{a-k}p\bn{m-k}p$ choices. Summing gives $U_N$.
A backward selected pair requires $k=p$, and must undo the whole
first hop, giving $W_p=L$. These are sums of squared CAR amplitudes;
their signs are not discarded by an absolute-value approximation.

For completeness the standard finite $G,G$ expression is
\begin{equation}\label{K}
 K_N=\frac{1}{\bn Nm}\sum_{r=0}^{\min(p,m-p)}
 \Lambda^r(m-p)\Lambda^r(p)
 \left[\bn Nr^2-\bn N{r-1}^2\right],
\end{equation}
with $\Lambda^r(k)=\bn{m-r}k\bn{N-m+k-r}k$; this is the
crossed term in the existing fourth-moment formula
\cite[Eq.~(28)]{Kota}. Alternatively $K_N$ can be left as the
raw CAR sum. Its double-scaled ratio $K_N/D_N^2\to q$ is the
ordinary complex DSSYK fourth-moment input already used in the
preceding stages \cite[Sec.~5.1]{BNR}. Unitary invariance makes the
annealed $H$-only Slater moments equal to the sector trace; the
half-filled sector has intersection parameter $e^{-4\lambda}$.
No finite identity inferred from a few fixtures
is needed here.

\section{Double-scaled mixed fourth moments}
Elementary binomial ratios and the preceding first-step deficit
law give
\begin{equation}\label{ratios}
 \frac LD\to s,\qquad r_N\to q,\qquad
 \frac{T_N}{D}\to q,\qquad\frac{U_N}{DL}\to q.
\end{equation}
For example, for fixed $c$,
\begin{equation}
 \log\frac{\bn{m-cp}p}{\bn mp}
 =\sum_{j=0}^{p-1}\log\left(1-\frac{cp}{m-j}\right)
 =-\frac{cp^2}{m}+o(1).
\end{equation}
For the last ratio in \eqref{ratios}, view $W_k/D$ as a probability.
The deficit $p-k$ converges to $\mathrm{Poisson}(2\lambda)$ from
the preceding stage. On any fixed deficit window the remaining
binomial product divided by $L$ converges uniformly to $q$.
That product lies in $[0,1]$, so tightness controls the window's
complement. This proves the limit without interchanging an
uncontrolled tail and the large-$N$ limit.
The nonzero backward terms vanish after normalization (as $1/L$
or $1/D$, with bounded mixed factors).

\newpage
One can now evaluate each limiting diagram with endpoint factors
$1$ for $HH$ or $MM^\dagger$, and $\gamma$ for $HM^\dagger$
or $MH$. Reject backward or incompatible selected edges.
Every crossing contributes $q$. A nesting contributes $q$
precisely when the inner edge is selected. All disjoint diagrams
and the remaining nesting have factor one. This is a
\emph{two-edge rule}; an arbitrary-edge version has not been proved.
Some useful moments are
\begin{align}
 \ph(h^4)&=2+q,&\ph(h^3M^\dagger)&=\gamma(2+q),\\
 \ph(h^2(M^\dagger)^2)&=2s q,
 &\ph(hM^\dagger hM^\dagger)&=s.
\end{align}
Here and below limiting notation abbreviates limits of the
specified finite sequences, not an assumed full representation.

\section{A positive space and the first Hamiltonian actions}
For polynomials $P,Q$ of degree at most two, set
$\langle P\Omega,Q\Omega\rangle=\lim_N\ph_N(P_N^\dagger Q_N)$.
Every entry exists by the table, and positivity follows from
finite-$N$ positivity. Quotient by null vectors. The identities
$M\Omega=0$, $MM^\dagger\Omega=\Omega$ in this quotient follow
from the preceding pure-ladder law. The additional relation is
\begin{equation}\label{Mh}
 M h\Omega=\gamma\Omega,
 \qquad
 \EE\|(M_Nh_N-\gamma_N I)\Om\|^2=\frac1{D_N}.
\end{equation}
Thus words ending in $M$ vanish on the vacuum, and the remaining
degree-two space is spanned by $I,M^\dagger,h,(M^\dagger)^2,
M^\dagger h,hM^\dagger,h^2$ acting on $\Omega$.

Define the following four centered degree-two vectors:
\begin{align}
 U&=(M^\dagger)^2\Omega,\\
 V&=\left(M^\dagger h-\gamma(M^\dagger)^2\right)\Omega,\\
 D_v&=\left(hM^\dagger-M^\dagger h-\gamma I\right)\Omega,\\
 R_v&=(h^2-I)\Omega-sU-2\gamma V-\gamma D_v.
\end{align}
They are mutually orthogonal and orthogonal to $\Omega$ and all
degree-one vectors. Their norms are
\begin{equation}\label{norms}
 \begin{split}
 \|U\|^2&=2q,\qquad \|V\|^2=q(1-s),\qquad\|D_v\|^2=1-q,\\
 \|R_v\|^2&=(1-s)^2(1+s+2q).
 \end{split}
\end{equation}
For example $\|hM^\dagger\Omega\|^2=1+s(1+q)$, while
$\|M^\dagger h\Omega\|^2=
\langle hM^\dagger\Omega,M^\dagger h\Omega\rangle=q(1+s)$;
these give the $D_v$ norm after subtracting its vacuum component.
All other entries follow from the same table. Gaussian parity
sets inner products with degree-one vectors to zero.

\newpage
Let $e_1=M^\dagger\Omega$, $e_r=(h\Omega-\gamma e_1)/\delta$,
and normalize the four vectors as
\begin{equation}
 e_2=\frac U{\sqrt{2q}},\quad
 e_{r2}=\frac V{\delta\sqrt q},\quad
 e_d=\frac{D_v}{\sqrt{1-q}},\quad
 e_{rr}=\frac{R_v}{(1-s)\sqrt{1+s+2q}}.
\end{equation}
For fixed $0<q<1$ every denominator is strictly positive.
The seven vectors $\Omega,e_1,e_r,e_2,e_{r2},e_d,e_{rr}$ are an
orthonormal basis of this degree-at-most-two quotient. Multiplication
by $h$ is determined from degree at most one into this space:
\begin{align}\label{actions}
 h\Omega&=\gamma e_1+\delta e_r,\\
 he_1&=\gamma\Omega+\gamma\sqrt{2q}\,e_2
       +\delta\sqrt q\,e_{r2}+\sqrt{1-q}\,e_d,\\
 he_r&=\delta\Omega+\gamma\sqrt q\,e_{r2}
       +\delta\sqrt{1+s+2q}\,e_{rr}.
\end{align}
These follow by substituting the four definitions; \eqref{Mh}
ensures the quotient relation needed below. No action of $h$ on
the degree-two basis is specified. Consequently this is not a
closed seven-state Hamiltonian or a finite-time solution.

The extra commutator direction has a particularly simple test:
\begin{equation}
 \|([h,M^\dagger]-\gamma I)\Omega\|^2=1-q>0.
\end{equation}
It is orthogonal to the pure second ladder state $e_2$, to the
mixed direction $e_{r2}$, and to the residual pair direction
$e_{rr}$. This resolves a definite part of the Hamiltonian leakage
instead of treating all missing norm as one scalar correction.

\section{The actual Gaussian-disorder complement}
At finite $N$ define
\begin{equation}
 g_N=h_N-\gamma_N(M_N+M_N^\dagger).
\end{equation}
This removes exactly the selected excitation block and its adjoint
from the original Hamiltonian. The remaining Gaussian coordinates
are independent of the selected coordinates. That statement is
about disorder coordinates; the CAR matrix factors still share
the same quantum state. The table gives
\begin{equation}\label{g}
 \ph(g^2)=1-s,\quad
 (g^2-(1-s)I)\Omega=R_v,\quad
 \frac{\ph(g^4)}{\ph(g^2)^2}=2+s+2q.
\end{equation}
To see the vector identity, expand $g^2\Omega$ using
$M\Omega=0$, $MM^\dagger\Omega=\Omega$, and \eqref{Mh}; the
result is exactly the definition of $R_v$ plus $(1-s)\Omega$.
Its norm then yields \eqref{g}. For $q>1/4$ this kurtosis exceeds
$3$. A standard normalized $q$-Gaussian mode would have fourth
moment $2+q$, so it cannot represent this complement.

\newpage
\section{Verification and the next calculation}
The independent finite CAR driver uses integer fermion actions
and Gaussian Wick sums, with no import of the old ladder driver.
It evaluates all 81 fourth words and all nine second words for
$(N,p)=(4,1),(6,1),(6,2),(8,2)$. In total it checks 972 pairing
entries from 2,000,499 Gaussian pair-label assignments, plus all
116 sector-basis columns of the channel acting on one disjoint
$p$-hop per fixture. It records six exact scalar ratio rows at
$\lambda_N=1/2$, $p=2,4,8,16,32,64$. Those asymptotic decimal
rows are diagnostics, not interval certificates.

The first driver enumerates the $HHHH$ crossed contribution; it
does not independently predict that contribution. A separate
read-only scalar check compares it with the literature expression
\eqref{K}. A second independent symbolic audit uses the limiting
two-edge rule to verify 62 identities, including the full $7\times7$
Gram matrix, \eqref{actions}, \eqref{g}, and four exact finite
residuals read from the CAR output. This checks algebraic
consequences, not a replacement for the analytic limiting proof.

The scientific records are
\begin{itemize}
 \item \texttt{results/susskind-slater-joint-fourth/exact-checks.json};
 \item \texttt{results/susskind-slater-joint-fourth/gram-checks.json};
 \item \texttt{results/susskind-slater-joint-fourth/covariance-checks.json}.
\end{itemize}
The stage evidence manifest binds source files, parameters,
environment versions, parent evidence, build logs, PDF and page
renders. Successful scientific outputs are preserved rather than
rerun for document compilation.

The next target is the all-degree same-$J$ joint law. At degree six,
three simultaneously active labels can test whether the two-edge
weights really factorize. A proof must retain general-hop overlaps,
CAR signs, and backward selected pairings; pure-ladder suppression
alone does not prove their suppression in the presence of $H$.
Once those rules are established, an all-layer positive
representation can supply the actual Hamiltonian evolution.
The first-layer calculation here does not yet identify a de Sitter
mass, position operator, stationary observer state, or finite-$G$
gravitational four-point function. Those are subsequent physics
questions, not substitutes for the joint-law calculation.

\begin{thebibliography}{9}
\bibitem{RX} H. Rajgadia and J. Xu,
\emph{Emergent States and Algebras from the Double-Scaling limit of
Pure States in SYK},
\href{https://arxiv.org/abs/2604.14387v3}{arXiv:2604.14387v3} (2026).
\bibitem{Kota} V. K. B. Kota,
\emph{SU(N) Wigner-Racah algebra for the matrix of second moments of
embedded Gaussian unitary ensemble of random matrices},
\href{https://arxiv.org/abs/nucl-th/0410069}{arXiv:nucl-th/0410069} (2004).
\bibitem{BNR} M. Berkooz, V. Narovlansky and H. Raj,
\emph{Complex Sachdev-Ye-Kitaev model in the double scaling limit},
\href{https://arxiv.org/abs/2006.13983v2}{arXiv:2006.13983v2} (2020).
\end{thebibliography}
\end{document}
